\documentclass[12pt]{article}

\usepackage[T2A]{fontenc}
\usepackage[utf8]{inputenc}
\usepackage[russian]{babel}
\usepackage{amsthm}
\usepackage{hyperref}
\usepackage{amsmath,amssymb}
\usepackage{booktabs}
\usepackage{longtable}
\usepackage{geometry}
\geometry{margin=1in}

\title{Связь таблицы скоростей продольных волн с уравнениями связанной термоупругости}
\author{}
\date{}

\begin{document}
\maketitle

\section*{1. Основная идея}

Таблица экспериментальных скоростей продольных волн $V_p(P)$ при различных всесторонних давлениях $P$ напрямую связана с комбинацией упругих модулей

\begin{equation}
V_p = \sqrt{\frac{\lambda + 2\mu}{\rho}}.
\end{equation}

Отсюда:
\begin{equation}
\lambda + 2\mu = \rho(P)\,V_p(P)^2.
\end{equation}

Таким образом, данные таблицы $V_p(P)$ подставляются в уравнения связанной термоупругости через величину
\begin{equation}
\lambda+2\mu.
\end{equation}
\section*{2. Разделение модулей}

Если имеется таблица скоростей поперечных волн $V_s(P)$, то
\begin{equation}
\mu(P) = \rho(P)\,V_s(P)^2,
\end{equation}
\begin{equation}
\lambda(P) = \rho(P)\,V_p(P)^2 - 2\mu(P),
\end{equation}
\begin{equation}
K(P) = \lambda(P) + \frac{2}{3}\mu(P).
\end{equation}

Если $V_s(P)$ нет, модуль $\mu$ можно оценить по литературным данным (числу Пуассона) или эмпирическим корреляциям.

\section*{3. Уравнение движения в термоупругости}

Для продольной волны в одномерном виде уравнение движения имеет вид:

\begin{equation}
\rho\,\frac{\partial^2 u}{\partial t^2}
=
(\lambda+2\mu)\frac{\partial^2 u}{\partial x^2}
-
3\lambda\alpha\,\frac{\partial T}{\partial x}.
\end{equation}

Именно в это уравнение подставляется величина $\lambda+2\mu = \rho V_p^2$.

\section*{4. Уравнение теплопроводности с учётом деформации}

\begin{equation}
k\,\frac{\partial^2 T}{\partial x^2}
=
\rho c\,\frac{\partial T}{\partial t}
+
3K\alpha T_0 \frac{\partial}{\partial t}\left(\frac{\partial u}{\partial x}\right).
\end{equation}

Где:
\begin{itemize}
    \item $k$ — теплопроводность,
    \item $c$ — удельная теплоёмкость,
    \item $T_0$ — температура образца,
    \item $\alpha$ — коэффициент линейного теплового расширения.
\end{itemize}

\section*{5. Связанная система в гармонической форме}

Для волновых решений вида $e^{i(kx-\omega t)}$ система принимает вид:
\begin{equation}
\begin{pmatrix}
(\lambda+2\mu)k^2 - \rho\omega^2 & -3\lambda\alpha\,k \\
i\omega\,3K\alpha T_0\,k & k\,k^2 + i\omega\rho c
\end{pmatrix}
\begin{pmatrix}
U\\
\Theta
\end{pmatrix}
= 0.
\end{equation}

Условие ненулевого решения:
\begin{equation}
\det(\cdot)=0,
\end{equation}
что даёт две продольные волны:
быструю (упругую) и медленную (термальную).

\section*{6. Какие таблицы необходимы}

Для корректного использования уравнений нужны следующие данные:

\begin{itemize}
    \item Таблица $V_p(P)$ — основная. (есть)
    
    \item Таблица $V_s(P)$ — для получения $\mu$. (есть)

    \item Таблица плотности $\rho(P)$. (есть)
    
    \item Коэффициент теплового расширения $\alpha(P,T)$. (есть)
    
    \item Удельная теплоёмкость $c(P,T)$. (есть)
    
    \item Теплопроводность $k(P,T)$. (есть)
    
    \item Температура образца $T_0$. (входные данные)
    
    \item (При пористости) пористость $\varphi(P)$ и свойства порового флюида. 
\end{itemize}

\section*{7. Итог}

Таблица $V_p(P)$ входит в уравнения термоупругости через комбинацию
\begin{equation}
\lambda + 2\mu = \rho(P)\,V_p(P)^2,
\end{equation}
что позволяет вычислить упругие параметры, необходимые для решения связанных уравнений и анализа быстрой и медленной продольной волны.

\section{Уравнения движения частицы, теплопроводности и связанной термоупругости}

Уравнение движения частицы:
\begin{equation}
\rho \frac{\partial^2 u_i}{\partial t^2} = \sigma_{ij,j} + f_i,
\end{equation}
где $u_i$ --- компоненты вектора перемещений, $\sigma_{ij}$ --- тензор напряжений, $f_i$ --- массовые силы.

Уравнение теплопроводности:
\begin{equation}
\rho c \frac{\partial T}{\partial t} = k \, \nabla^2 T + Q,
\end{equation}
где $T$ — температура, $k$ — коэффициент теплопроводности, $Q$ — источники тепла.

Связанная система термоупругости объединяет оба уравнения с учётом теплового расширения:
\begin{align}
\rho \frac{\partial^2 u_i}{\partial t^2} &= C_{ijkl} u_{k,lj} - 3 K \alpha \, T_{,i}, \\
\rho c \frac{\partial T}{\partial t} &= k \, T_{,ii} + 3 K \alpha T_0 \, u_{k,k t},
\end{align}
где $K$ — модуль объёмного сжатия, $\alpha$ — коэффициент линейного теплового расширения, $T_0$ — начальная температура.

\section{Таблица свойств горных пород при высоких давлениях}

\begin{longtable}{lccccccccc}
\caption{Изменение величины отношения $k = V_p / V_s$ при высоких давлениях} \\ 
\toprule
Порода & Параметр & 1 & 500 & 1000 & 2000 & 4000 & 6000 & 10000 & 15000 \\
\midrule
Эклогит 91-52 & $V_p$ & 6.17 & 7.04 & 7.23 & 7.42 & 7.59 & 7.79 & 7.89 & 8.00 \\
 & $V_s$ & 3.77 & 3.86 & 3.87 & 3.92 & 3.97 & 4.00 & 4.06 & 4.10 \\
 & $k$ & 1.64 & 1.82 & 1.86 & 1.89 & 1.91 & 1.95 & 1.95 & 1.95 \\
\midrule
Вебстерит плагиоклазовый 86-25(1) & $V_p$ & 6.11 & 6.44 & 6.46 & 6.48 & 6.54 & 6.60 & 6.74 & 6.90 \\
 & $V_s$ & 3.91 & 4.01 & 4.07 & 4.08 & 4.08 & 4.10 & 4.11 & 4.12 \\
 & $k$ & 1.56 & 1.61 & 1.59 & 1.59 & 1.60 & 1.61 & 1.64 & 1.68 \\
\midrule
Кварц-биотитовая порода 88-2(3) & $V_p$ & 6.12 & 6.80 & 6.90 & 7.08 & 7.29 & 7.40 & 7.58 & 7.82 \\
 & $V_s$ & 3.52 & 3.58 & 3.64 & 3.72 & 3.80 & 3.82 & 3.84 & 3.88 \\
 & $k$ & 1.74 & 1.85 & 1.89 & 1.90 & 1.92 & 1.94 & 1.97 & 2.02 \\
\midrule
Актиноилит 209-79(1) & $V_p$ & 5.38 & 5.90 & 6.07 & 6.26 & 6.60 & 6.82 & 6.94 & 7.10 \\
 & $V_s$ & 2.74 & 3.10 & 3.50 & 3.66 & 3.69 & 3.72 & 3.78 & 3.80 \\
 & $k$ & 1.96 & 1.90 & 1.73 & 1.71 & 1.79 & 1.83 & 1.84 & 1.86 \\
\midrule
Амфиболит гранатовый 11-12(3) & $V_p$ & 6.24 & 7.00 & 7.06 & 7.23 & 7.37 & 7.46 & 7.63 & 7.82 \\
 & $V_s$ & 3.80 & 3.84 & 3.89 & 3.94 & 3.99 & 4.05 & 4.08 & 4.12 \\
 & $k$ & 1.57 & 1.82 & 1.81 & 1.83 & 1.85 & 1.84 & 1.87 & 1.90 \\
\bottomrule
\end{longtable}

\end{document}
